\section{Busqueda intra-sistema de candidatos a planetas no detectados}

\subsection{Motivacion}

Este modulo traslada la incompletitud topologica global a la escala de sistemas planetarios conocidos. El objetivo no es afirmar detecciones fisicas directas, sino priorizar intervalos orbitales donde podria existir submuestreo intra-sistema.

\subsection{Metodologia}

Cada sistema se ordena por periodo orbital y se analizan gaps adyacentes. Para gaps grandes se generan candidatos sinteticos en escala logaritmica de periodo y semieje mayor, con apoyo adicional de analogos globales y proxies de detectabilidad por transito o velocidad radial.

\subsection{Uso de resultados previos TOI/ATI/sombra}

Cuando existen tablas previas, el modulo incorpora TOI regional, ATI por planeta ancla y sombra observacional por nodo Mapper. Esto convierte la topologia global mas la sombra observacional mas la arquitectura intra-sistema en un ranking de intervalos orbitales prioritarios.

\subsection{Resultados principales}

Los resultados priorizan candidatos topologicos y regiones sugeridas para seguimiento observacional. No se detectan planetas reales ni se afirma un conteo verdadero de planetas ausentes.

\begin{itemize}

\item HIP 90988: max score=0.000, n candidatos=0.

\end{itemize}

\subsection{Validacion leave-one-out}

La validacion interna retira temporalmente planetas intermedios conocidos y verifica si el modulo vuelve a priorizar candidatos cercanos en periodo orbital.

Se ejecutaron 0 pruebas. El error mediano en log10(P) fue nan, con recall dentro de 0.1 dex de nan y dentro de 0.2 dex de nan.

\subsection{Limitaciones}

El ranking no modela completitud instrumental ni sustituye analisis de inyeccion-recuperacion. Las estimaciones de masa, radio y detectabilidad son proxies para priorizacion, no inferencias fisicas definitivas.

\subsection{Trabajo futuro}

Los siguientes pasos incluyen incorporar completitud instrumental, seguimiento observacional real y validacion cruzada con catalogos futuros.
